\documentclass{article}%
\usepackage[T1]{fontenc}%
\usepackage[utf8]{inputenc}%
\usepackage{lmodern}%
\usepackage{textcomp}%
\usepackage{lastpage}%
\usepackage{geometry}%
\geometry{margin=1in,headheight=14pt,headsep=25pt}%
\usepackage{hyperref}%
\usepackage{enumitem}%
\usepackage{fancyhdr}%
\usepackage{xcolor}%
\usepackage{url}%
\usepackage{breakurl}%
%

            \hypersetup{
                colorlinks=true,
                linkcolor=blue,
                filecolor=magenta,
                urlcolor=blue
            }
            \pagestyle{fancy}
            \fancyhf{}
            \rhead{Generated on \today}
            \cfoot{\thepage}
            
            % Configure enumeration settings
            \setlist[enumerate,1]{label=\arabic*.}
            \setlist[enumerate,2]{label=\alph*.}
            \setlist[enumerate,3]{label=\Alph*.}
            \setlist[enumerate]{itemsep=0.5em}
            
            % Define a command for red text
            \newcommand{\incorrect}[1]{\textcolor{red}{#1}}
        %
\lhead{2024{-}09{-}17{-}DualityGeneral}%
\title{Practice Questions for 2024{-}09{-}17{-}DualityGeneral}%
\author{Generated by Scribe.AI}%
\date{\today}%
%
\begin{document}%
\normalsize%
\maketitle%
\section{Questions}%
\label{sec:Questions}%
\begin{enumerate}%
\item What is the significance of the Lagrangian function in deriving the dual problem from a primal linear program with both equality and inequality constraints?%
\begin{enumerate}%
\item The Lagrangian function directly provides the objective function of the dual problem.%
\item The Lagrangian function transforms the constrained primal problem into an unconstrained problem, facilitating the derivation of the dual problem.%
\item The Lagrangian function is only relevant for primal problems with equality constraints, not inequality constraints.%
\item The Lagrangian function is irrelevant to the derivation of the dual problem; the dual problem is derived solely from the primal problem's constraints.%
\item The Lagrangian function provides the constraints of the dual problem.%
\end{enumerate}%
\vspace{1em}%
\item In the context of weak duality, what is the relationship between the optimal objective function values of the primal and dual linear programs?%
\begin{enumerate}%
\item The optimal objective function value of the primal problem is always greater than or equal to the optimal objective function value of the dual problem.%
\item The optimal objective function values of the primal and dual problems are always equal.%
\item The optimal objective function value of the primal problem is always less than the optimal objective function value of the dual problem.%
\item There is no relationship between the optimal objective function values of the primal and dual problems.%
\item The optimal objective function value of the primal problem is always less than or equal to the optimal objective function value of the dual problem.%
\end{enumerate}%
\vspace{1em}%
\item What condition must be satisfied for the dual objective function to be finite (not $+\infty$) when deriving the dual from a general linear program using the Lagrangian method?%
\begin{enumerate}%
\item The primal problem must be infeasible.%
\item The vector of primal variables must be zero.%
\item The Lagrange multipliers must satisfy a specific linear equation involving the constraint matrices and the primal objective function vector.%
\item The dual problem must be infeasible.%
\item The primal objective function must be zero.%
\end{enumerate}%
\vspace{1em}%
\item How does the non-negativity constraint on the Lagrange multipliers associated with inequality constraints in the primal problem affect the derivation of the dual problem?%
\begin{enumerate}%
\item It has no effect; the non-negativity constraint is only relevant for the primal problem.%
\item It leads to equality constraints in the dual problem.%
\item It leads to inequality constraints in the dual problem.%
\item It ensures the dual objective function is always bounded.%
\item It is not considered during the derivation of the dual problem.%
\end{enumerate}%
\vspace{1em}%
\item Explain the concept of weak duality in the context of primal and dual linear programs, and how it relates to the feasibility and boundedness of both problems.%
\begin{enumerate}%
\item Weak duality states that the optimal objective function value of the primal problem is always less than or equal to the optimal objective function value of the dual problem, regardless of feasibility.%
\item Weak duality only applies if both the primal and dual problems are feasible and bounded.  If either is infeasible or unbounded, weak duality does not hold.%
\item Weak duality implies that if the primal problem is unbounded, then the dual problem must be infeasible, and vice versa.%
\item Weak duality states that the optimal objective function value of the primal problem is always greater than or equal to the optimal objective function value of the dual problem, regardless of feasibility.%
\item Weak duality is only relevant for linear programs with equality constraints; it does not apply to linear programs with inequality constraints.%
\end{enumerate}%
\vspace{1em}%
\item Given the primal problem $\max \quad \xi(x) = c^T x$ subject to $Ax \le b, x \ge 0$, and its dual problem $\min \quad \xi(y) = b^T y$ subject to $A^T y \ge c, y \ge 0$, what is the value of $\xi(y)$ if $b = \begin{pmatrix} 1 \\ 2 \end{pmatrix}$ and $y = \begin{pmatrix} 1 \\ 1 \end{pmatrix}$?%
\begin{enumerate}%
\item 1%
\item 2%
\item 3%
\item 0%
\item 4%
\end{enumerate}%
\vspace{1em}%
\item Consider the primal problem: $\max \quad \xi(x) = 2x_1 + 3x_2$ subject to $x_1 + x_2 \le 4$, $2x_1 + x_2 \le 5$, $x_1, x_2 \ge 0$. What is the second constraint of the dual problem?%
\begin{enumerate}%
\item $y_1 + y_2 \ge 3$%
\item $y_1 + 2y_2 \ge 2$%
\item $2y_1 + y_2 \ge 3$%
\item $4y_1 + 5y_2 \ge 2$%
\item $y_1 + y_2 \ge 2$%
\end{enumerate}%
\vspace{1em}%
\item Given the primal problem $\max \quad \xi(x) = c^T x$ subject to $Ax \le b, x \ge 0$, and its dual problem $\min \quad \xi(y) = b^T y$ subject to $A^T y \ge c, y \ge 0$, if $A = \begin{pmatrix} 1 & 2 \\ 3 & 4 \end{pmatrix}$, $b = \begin{pmatrix} 5 \\ 6 \end{pmatrix}$, and $c = \begin{pmatrix} 1 \\ 1 \end{pmatrix}$, what is the dual objective function $\xi(y)$?%
\begin{enumerate}%
\item $\xi(y) = y_1 + y_2$%
\item $\xi(y) = 5y_1 + 6y_2$%
\item $\xi(y) = y_1 + 2y_2$%
\item $\xi(y) = 6y_1 + 5y_2$%
\item $\xi(y) = 3y_1 + 4y_2$%
\end{enumerate}%
\vspace{1em}%
\item Write the dual problem (D) explicitly, given the primal problem (P): $\max \quad \xi(x) = x_1 + 2x_2$ subject to $x_1 + x_2 \le 4$, $2x_1 + x_2 \le 5$, $x_1, x_2 \ge 0$.%
\begin{enumerate}%
\item $\min \quad \xi(y) = 4y_1 + 5y_2$ subject to $y_1 + 2y_2 \ge 1$, $y_1 + y_2 \ge 2$, $y_1, y_2 \ge 0$%
\item $\min \quad \xi(y) = 5y_1 + 4y_2$ subject to $y_1 + 2y_2 \ge 1$, $y_1 + y_2 \ge 2$, $y_1, y_2 \ge 0$%
\item $\min \quad \xi(y) = 4y_1 + 5y_2$ subject to $y_1 + 2y_2 \le 1$, $y_1 + y_2 \le 2$, $y_1, y_2 \ge 0$%
\item $\max \quad \xi(y) = 4y_1 + 5y_2$ subject to $y_1 + 2y_2 \ge 1$, $y_1 + y_2 \ge 2$, $y_1, y_2 \ge 0$%
\item $\min \quad \xi(y) = 4y_1 + 5y_2$ subject to $2y_1 + y_2 \ge 1$, $y_1 + 2y_2 \ge 2$, $y_1, y_2 \ge 0$%
\end{enumerate}%
\vspace{1em}%
\item%
\begin{enumerate}%
\item Write the dual problem (D) explicitly, substituting the given values of A, b, and c.%
\begin{enumerate}%
\item Minimize $\xi(y) = 5y_1 + 6y_2$ subject to $y_1 + 3y_2 \ge 1$, $2y_1 + 4y_2 \ge 1$, $y_1, y_2 \ge 0$%
\item Maximize $\xi(y) = 5y_1 + 6y_2$ subject to $y_1 + 3y_2 \le 1$, $2y_1 + 4y_2 \le 1$, $y_1, y_2 \ge 0$%
\item Minimize $\xi(y) = 5y_1 + 6y_2$ subject to $y_1 + 3y_2 \le 1$, $2y_1 + 4y_2 \le 1$, $y_1, y_2 \ge 0$%
\item Maximize $\xi(y) = 5y_1 + 6y_2$ subject to $y_1 + 3y_2 \ge 1$, $2y_1 + 4y_2 \ge 1$, $y_1, y_2 \ge 0$%
\item Minimize $\xi(y) = y_1 + y_2$ subject to $y_1 + 2y_2 \le 5$, $3y_1 + 4y_2 \le 6$, $y_1, y_2 \ge 0$%
\end{enumerate}%
\vspace{0.5em}%
\item If $x = \begin{pmatrix} 0 \\ 2 \end{pmatrix}$ is a feasible solution for the primal problem (P), what is the value of the primal objective function $\xi(x)$?%
\begin{enumerate}%
\item 0%
\item 2%
\item 5%
\item 6%
\item 10%
\end{enumerate}%
\vspace{0.5em}%
\item Is the solution $x = \begin{pmatrix} 0 \\ 2 \end{pmatrix}$ feasible for the primal problem (P)? Show your work.%
\begin{enumerate}%
\item Yes%
\item No%
\item Cannot be determined%
\item Only if $c = \begin{pmatrix} 0 \\ 1 \end{pmatrix}$%
\item Only if $b = \begin{pmatrix} 4 \\ 8 \end{pmatrix}$%
\end{enumerate}%
\vspace{0.5em}%
\item Suppose a feasible solution to the dual problem is $y = \begin{pmatrix} 1/2 \\ 1/2 \end{pmatrix}$. What is the value of the dual objective function $\xi(y)$?%
\begin{enumerate}%
\item 1/2%
\item 1%
\item 5.5%
\item 6.5%
\item 11/2%
\end{enumerate}%
\vspace{0.5em}%
\item Is the solution $y = \begin{pmatrix} 1/2 \\ 1/2 \end{pmatrix}$ feasible for the dual problem (D)? Show your work.%
\begin{enumerate}%
\item Yes%
\item No%
\item Cannot be determined%
\item Only if $c = \begin{pmatrix} 0 \\ 1 \end{pmatrix}$%
\item Only if $b = \begin{pmatrix} 4 \\ 8 \end{pmatrix}$%
\end{enumerate}%
\vspace{0.5em}%
\end{enumerate}%
\end{enumerate}

%
\newpage%
\section{Answers}%
\label{sec:Answers}%
\begin{enumerate}%
\item What is the significance of the Lagrangian function in deriving the dual problem from a primal linear program with both equality and inequality constraints?%
\begin{enumerate}%
\item \incorrect{The Lagrangian function does not directly give the dual objective function.  The dual objective function is derived from the Lagrangian after maximizing it with respect to the primal variables.}%
\item The Lagrangian function is a crucial tool in deriving the dual problem. By incorporating the constraints into the objective function using Lagrange multipliers, it allows us to analyze the problem without explicitly considering the constraints, simplifying the derivation of the dual.%
\item \incorrect{The Lagrangian function is essential for handling both equality and inequality constraints in the primal problem.  The Lagrange multipliers associated with inequality constraints are non-negative.}%
\item \incorrect{The Lagrangian function is fundamental to the derivation of the dual problem.  It provides a way to incorporate the constraints into the objective function, which is essential for finding the dual.}%
\item \incorrect{The Lagrangian function does not directly provide the constraints of the dual problem. The constraints of the dual problem are derived from the conditions that ensure the Lagrangian is bounded when maximizing with respect to the primal variables.}%
\end{enumerate}%
\vspace{1em}%
\item In the context of weak duality, what is the relationship between the optimal objective function values of the primal and dual linear programs?%
\begin{enumerate}%
\item \incorrect{This is incorrect.  It describes strong duality, which only holds under certain conditions. Weak duality states that the primal optimal value is less than or equal to the dual optimal value.}%
\item \incorrect{This is incorrect. This describes strong duality, which only holds under certain conditions. Weak duality only guarantees an inequality between the optimal values.}%
\item \incorrect{This is incorrect. It is the opposite of weak duality. Weak duality states that the primal optimal value is less than or equal to the dual optimal value.}%
\item \incorrect{This is incorrect. Weak duality establishes a fundamental relationship between the optimal values of the primal and dual problems.}%
\item Weak duality states that the optimal objective function value of the primal problem is always less than or equal to the optimal objective function value of the dual problem. This holds regardless of whether the problems are feasible or infeasible.%
\end{enumerate}%
\vspace{1em}%
\item What condition must be satisfied for the dual objective function to be finite (not $+\infty$) when deriving the dual from a general linear program using the Lagrangian method?%
\begin{enumerate}%
\item \incorrect{The feasibility of the primal problem does not directly determine the finiteness of the dual objective function.  The dual objective function can be finite even if the primal is infeasible (though the dual would be unbounded in that case).}%
\item \incorrect{The value of the primal variables is irrelevant to the finiteness of the dual objective function. The dual objective function is derived by maximizing the Lagrangian with respect to the primal variables.}%
\item The dual objective function is finite only if the coefficient of the primal variables in the Lagrangian is zero. This condition translates into a linear equation involving the Lagrange multipliers, the constraint matrices from the primal problem, and the primal objective function vector.%
\item \incorrect{The feasibility of the dual problem is related to the boundedness of the dual objective function, but it doesn't directly determine whether the dual objective function is finite or infinite.  The condition for a finite dual objective function is a constraint on the Lagrange multipliers.}%
\item \incorrect{The value of the primal objective function is not directly related to the finiteness of the dual objective function. The condition for a finite dual objective function is a constraint on the Lagrange multipliers.}%
\end{enumerate}%
\vspace{1em}%
\item How does the non-negativity constraint on the Lagrange multipliers associated with inequality constraints in the primal problem affect the derivation of the dual problem?%
\begin{enumerate}%
\item \incorrect{The non-negativity constraint on the Lagrange multipliers is crucial for the derivation of the dual problem. It directly impacts the constraints of the dual problem.}%
\item \incorrect{The non-negativity constraint on the Lagrange multipliers leads to inequality constraints, not equality constraints, in the dual problem.}%
\item The non-negativity constraint on the Lagrange multipliers associated with inequality constraints in the primal problem translates directly into inequality constraints in the dual problem. This is a fundamental aspect of duality theory.%
\item \incorrect{While the non-negativity constraint helps ensure the dual problem is well-behaved, it doesn't guarantee that the dual objective function is always bounded. The boundedness depends on the feasibility of the primal and dual problems.}%
\item \incorrect{The non-negativity constraint is explicitly considered during the derivation of the dual problem. It is a key element in transforming the primal problem into its dual.}%
\end{enumerate}%
\vspace{1em}%
\item Explain the concept of weak duality in the context of primal and dual linear programs, and how it relates to the feasibility and boundedness of both problems.%
\begin{enumerate}%
\item Weak duality states that the optimal objective function value of the primal problem is always less than or equal to the optimal objective function value of the dual problem. This holds even if one or both problems are infeasible.  If the primal is unbounded, the dual is infeasible, and vice versa. This relationship provides valuable information about the primal and dual problems' solutions, even before finding the optimal solutions.%
\item \incorrect{This statement is incorrect. Weak duality holds regardless of the feasibility or boundedness of the primal and dual problems.  If the primal is infeasible, the dual may be infeasible or unbounded. If the primal is unbounded, the dual is infeasible.  If the dual is infeasible, the primal may be infeasible or unbounded. If the dual is unbounded, the primal is infeasible.}%
\item \incorrect{This statement is partially correct in that if the primal is unbounded, the dual is infeasible, and vice versa. However, it doesn't encompass the full meaning of weak duality, which also applies when both problems are feasible and bounded.}%
\item \incorrect{This statement is incorrect. Weak duality states that the optimal objective function value of the primal problem is less than or equal to the optimal objective function value of the dual problem.}%
\item \incorrect{This statement is incorrect. Weak duality applies to both linear programs with equality and inequality constraints.}%
\end{enumerate}%
\vspace{1em}%
\item Given the primal problem $\max \quad \xi(x) = c^T x$ subject to $Ax \le b, x \ge 0$, and its dual problem $\min \quad \xi(y) = b^T y$ subject to $A^T y \ge c, y \ge 0$, what is the value of $\xi(y)$ if $b = \begin{pmatrix} 1 \\ 2 \end{pmatrix}$ and $y = \begin{pmatrix} 1 \\ 1 \end{pmatrix}$?%
\begin{enumerate}%
\item \incorrect{This would be the result if $b = \begin{pmatrix} 1 \\ 0 \end{pmatrix}$ and $y = \begin{pmatrix} 1 \\ 1 \end{pmatrix}$}%
\item \incorrect{This would be the result if $b = \begin{pmatrix} 0 \\ 2 \end{pmatrix}$ and $y = \begin{pmatrix} 1 \\ 1 \end{pmatrix}$}%
\item The dual objective function is given by $\xi(y) = b^T y$. Substituting the given values, we have $\xi(y) = \begin{pmatrix} 1 & 2 \end{pmatrix} \begin{pmatrix} 1 \\ 1 \end{pmatrix} = 1 + 2 = 3$. Therefore, the value of $\xi(y)$ is 3.%
\item \incorrect{This would be the result if $b = \begin{pmatrix} 0 \\ 0 \end{pmatrix}$ and $y = \begin{pmatrix} 1 \\ 1 \end{pmatrix}$}%
\item \incorrect{This would be the result if $b = \begin{pmatrix} 2 \\ 2 \end{pmatrix}$ and $y = \begin{pmatrix} 1 \\ 1 \end{pmatrix}$}%
\end{enumerate}%
\vspace{1em}%
\item Consider the primal problem: $\max \quad \xi(x) = 2x_1 + 3x_2$ subject to $x_1 + x_2 \le 4$, $2x_1 + x_2 \le 5$, $x_1, x_2 \ge 0$. What is the second constraint of the dual problem?%
\begin{enumerate}%
\item \incorrect{This mixes up the coefficients from the primal objective function and constraints.}%
\item \incorrect{This is the first constraint of the dual problem.}%
\item The dual problem is $\min \quad \xi(y) = 4y_1 + 5y_2$ subject to $y_1 + 2y_2 \ge 2$, $y_1 + y_2 \ge 3$, $y_1, y_2 \ge 0$. The second constraint comes from the coefficient of $x_2$ in the primal constraints.  Therefore, the second constraint of the dual problem is $y_1 + y_2 \ge 3$.%
\item \incorrect{This incorrectly combines the right-hand side values of the primal constraints.}%
\item \incorrect{This is incorrect; it uses the correct variables but the wrong coefficients.}%
\end{enumerate}%
\vspace{1em}%
\item Given the primal problem $\max \quad \xi(x) = c^T x$ subject to $Ax \le b, x \ge 0$, and its dual problem $\min \quad \xi(y) = b^T y$ subject to $A^T y \ge c, y \ge 0$, if $A = \begin{pmatrix} 1 & 2 \\ 3 & 4 \end{pmatrix}$, $b = \begin{pmatrix} 5 \\ 6 \end{pmatrix}$, and $c = \begin{pmatrix} 1 \\ 1 \end{pmatrix}$, what is the dual objective function $\xi(y)$?%
\begin{enumerate}%
\item \incorrect{This uses the values from $c$ instead of $b$.}%
\item The dual objective function is given by $\xi(y) = b^T y$.  Substituting the given value of $b$, we get $\xi(y) = \begin{pmatrix} 5 & 6 \end{pmatrix} \begin{pmatrix} y_1 \\ y_2 \end{pmatrix} = 5y_1 + 6y_2$.%
\item \incorrect{This uses some elements from $A$ and $b$ incorrectly.}%
\item \incorrect{This swaps the elements of $b$.}%
\item \incorrect{This uses elements from $A$ instead of $b$.}%
\end{enumerate}%
\vspace{1em}%
\item Write the dual problem (D) explicitly, given the primal problem (P): $\max \quad \xi(x) = x_1 + 2x_2$ subject to $x_1 + x_2 \le 4$, $2x_1 + x_2 \le 5$, $x_1, x_2 \ge 0$.%
\begin{enumerate}%
\item The dual problem is formed by taking the transpose of the constraint matrix A, and using the vector b for the objective function. The inequalities are reversed, and the objective function is minimized. Thus, the dual problem is $\min \quad \xi(y) = 4y_1 + 5y_2$ subject to $y_1 + 2y_2 \ge 1$, $y_1 + y_2 \ge 2$, $y_1, y_2 \ge 0$.%
\item \incorrect{The objective function coefficients are incorrect.}%
\item \incorrect{The direction of the inequalities is incorrect.}%
\item \incorrect{The objective function is maximized instead of minimized.}%
\item \incorrect{The constraints are incorrectly formed.}%
\end{enumerate}%
\vspace{1em}%
\item%
\begin{enumerate}%
\item Write the dual problem (D) explicitly, substituting the given values of A, b, and c.%
\begin{enumerate}%
\item The dual problem is formed by taking the transpose of A, switching the objective function from maximization to minimization, and changing the direction of the inequalities. The objective function becomes $b^T y = 5y_1 + 6y_2$, and the constraints become $A^T y \ge c$, which translates to $y_1 + 3y_2 \ge 1$ and $2y_1 + 4y_2 \ge 1$. The non-negativity constraints $y_1, y_2 \ge 0$ remain.%
\item \incorrect{This option incorrectly changes the direction of the inequalities in the constraints.}%
\item \incorrect{This option incorrectly changes the direction of the inequalities and the type of the objective function.}%
\item \incorrect{This option incorrectly changes the type of the objective function.}%
\item \incorrect{This option incorrectly uses $c$ instead of $b$ in the objective function and incorrectly reverses the inequalities in the constraints.}%
\end{enumerate}%
\vspace{0.5em}%
\item If $x = \begin{pmatrix} 0 \\ 2 \end{pmatrix}$ is a feasible solution for the primal problem (P), what is the value of the primal objective function $\xi(x)$?%
\begin{enumerate}%
\item \incorrect{This is incorrect; it assumes $x_2$ is 0.}%
\item Substituting $x = \begin{pmatrix} 0 \\ 2 \end{pmatrix}$ into $\xi(x) = c^T x = \begin{pmatrix} 1 & 1 \end{pmatrix} \begin{pmatrix} 0 \\ 2 \end{pmatrix} = 2$.%
\item \incorrect{This is incorrect; it uses a component of $b$ instead of $c$.}%
\item \incorrect{This is incorrect; it uses a component of $b$ instead of $c$.}%
\item \incorrect{This is incorrect; it multiplies the components of $x$ and $c$ incorrectly.}%
\end{enumerate}%
\vspace{0.5em}%
\item Is the solution $x = \begin{pmatrix} 0 \\ 2 \end{pmatrix}$ feasible for the primal problem (P)? Show your work.%
\begin{enumerate}%
\item We check the constraints: $Ax = \begin{pmatrix} 1 & 2 \\ 3 & 4 \end{pmatrix} \begin{pmatrix} 0 \\ 2 \end{pmatrix} = \begin{pmatrix} 4 \\ 8 \end{pmatrix}$. Since $\begin{pmatrix} 4 \\ 8 \end{pmatrix} \le \begin{pmatrix} 5 \\ 6 \end{pmatrix}$ and $x \ge 0$, the solution is feasible.%
\item \incorrect{This is incorrect; the solution satisfies the constraints.}%
\item \incorrect{This is incorrect; feasibility can be determined by checking the constraints.}%
\item \incorrect{This is incorrect; the feasibility of $x$ is independent of $c$.}%
\item \incorrect{This is incorrect; the feasibility of $x$ is independent of $b$.}%
\end{enumerate}%
\vspace{0.5em}%
\item Suppose a feasible solution to the dual problem is $y = \begin{pmatrix} 1/2 \\ 1/2 \end{pmatrix}$. What is the value of the dual objective function $\xi(y)$?%
\begin{enumerate}%
\item \incorrect{This is incorrect; it is half the correct answer.}%
\item \incorrect{This is incorrect; it is half the correct answer.}%
\item Substituting $y = \begin{pmatrix} 1/2 \\ 1/2 \end{pmatrix}$ into $\xi(y) = b^T y = \begin{pmatrix} 5 & 6 \end{pmatrix} \begin{pmatrix} 1/2 \\ 1/2 \end{pmatrix} = 5/2 + 6/2 = 11/2 = 5.5$.%
\item \incorrect{This is incorrect; it is close but not the correct answer.}%
\item \incorrect{This is the correct answer in fraction form.}%
\end{enumerate}%
\vspace{0.5em}%
\item Is the solution $y = \begin{pmatrix} 1/2 \\ 1/2 \end{pmatrix}$ feasible for the dual problem (D)? Show your work.%
\begin{enumerate}%
\item \incorrect{This is incorrect; the solution does not satisfy the constraints.}%
\item We check the constraints: $A^T y = \begin{pmatrix} 1 & 3 \\ 2 & 4 \end{pmatrix} \begin{pmatrix} 1/2 \\ 1/2 \end{pmatrix} = \begin{pmatrix} 2 \\ 3 \end{pmatrix}$. Since $\begin{pmatrix} 2 \\ 3 \end{pmatrix} \ge \begin{pmatrix} 1 \\ 1 \end{pmatrix}$ and $y \ge 0$, the solution is feasible.%
\item \incorrect{This is incorrect; feasibility can be determined by checking the constraints.}%
\item \incorrect{This is incorrect; the feasibility of $y$ is independent of $c$.}%
\item \incorrect{This is incorrect; the feasibility of $y$ is independent of $b$.}%
\end{enumerate}%
\vspace{0.5em}%
\end{enumerate}%
\end{enumerate}

%
\end{document}